\chapter{Twierdzenie Smale'a o \texorpdfstring{$h$}{h}-kobordyzmie}
\label{ch:hcobordyzm}
\index{h-kobordyzm@$h$-kobordyzm}

W poprzednim rozdziale nauczyliśmy się
przesuwać uchwyty i znosić parę, której
sfery spotykają się w jednym punkcie.
Teraz zastosujemy te ruchy do całego
kobordyzmu. Intuicja jest prosta:
jeśli oba brzegi $W$ zachowują w nim
całą informację homotopijną, wszystkie
uchwyty powinny występować w parach,
które nie wnoszą trwałej topologii.
Dowód musi jeszcze zamienić to zdanie
o \emph{homotopii} na stwierdzenie
o \emph{dyfeomorfizmie}.

Przez cały rozdział $(W;N_-,N_+)$
oznacza zwarty, spójny, gładki
kobordyzm wymiaru $m$ z dwoma
niepustymi spójnymi brzegami.
Używamy oznaczeń rdzenia, sfery
przyczepienia $A$ i sfery pasa $B$
z Rozdziału~\ref{ch:rachunek-uchwytow}.
Liczby przecięcia dotyczą poziomicy
wymiaru $m-1$.

\section{Co dokładnie orzeka twierdzenie?}
\label{sec:stwierdzenie-hc-19}

\begin{twierdzenie}[Smale, twierdzenie o $h$-kobordyzmie]
\label{thm:smale-hc-19}
Załóżmy, że $m\geq6$, rozmaitość
$W^m$ jest po prostu spójna, a obie
inkluzje $N_\pm\hookrightarrow W$
są równoważnościami homotopijnymi.
Wówczas istnieje dyfeomorfizm
\[
 \Phi:N_-\times[0,1]\longrightarrow W,
 \qquad \Phi(x,0)=x.
\]
W szczególności $N_-\cong N_+$
jako rozmaitości gładkie.
\end{twierdzenie}

To twierdzenie dotyczy \emph{gładkiej}
struktury $W$; końcowy dyfeomorfizm
na $N_+$ nie jest z góry wybrany.
Termin \emph{$h$-kobordyzm} dla
warunku o obu inkluzjach pojawił się
w Sekcji~\ref{sec:uchwyty-smale-most-18}.
Prosta spójność $W$ wraz z tymi
równoważnościami oznacza też
$\pi_1(N_-)=\pi_1(N_+)=0$.

\begin{stwierdzenie}[Acykliczność]
\label{prop:acykliczny-hc-19}
W warunkach
Twierdzenia~\ref{thm:smale-hc-19}
zachodzi
\[
 H_j(W,N_-;\Z)=0
 \qquad\text{dla każdego }j.
\]
Kompleks uchwytów z
Twierdzenia~\ref{thm:kompleks-uchwytow-18}
jest zatem acykliczny.
\end{stwierdzenie}
\begin{proof}
Równoważność homotopijna
$N_-\hookrightarrow W$ indukuje
izomorfizm wszystkich grup homologii.
W długim ciągu dokładnym pary
\[
 \cdots\to H_j(N_-)\to H_j(W)
 \to H_j(W,N_-)\to H_{j-1}(N_-)
 \to H_{j-1}(W)\to\cdots
\]
obie widoczne mapy indukowane inkluzją
są izomorfizmami. Dokładność wymusza
zatem zerowanie grupy pośrodku.
Kompleks uchwytów oblicza te grupy,
więc jego homologia jest zerowa.
\end{proof}

Acykliczność mówi, że macierze
$\partial_k$ zawierają informację
o możliwych \emph{algebraicznych}
parach. Nie mówi jeszcze, czy
sfery spotykają się tylko raz.
Rysunek~\ref{fig:plan-smale-19}
pokazuje drogę od tego faktu
do iloczynu.

\begin{figure}[htbp]
\centering
\begin{tikzpicture}[>=Latex,font=\small,
 every node/.style={align=center,rounded corners=3pt,
                    inner sep=5pt}]
 \node[draw=manifoldblue!70!black,fill=manifoldblue!9]
      (a) at (-2.55,.90) {acykliczny kompleks\\uchwytów};
 \node[draw=manifoldgold!85!black,fill=manifoldgold!12]
      (b) at (2.55,.90) {przesunięcia\\macierz z blokami $1$};
 \node[draw=manifoldgreen!70!black,fill=manifoldgreen!10]
      (c) at (2.55,-.90) {trik Whitneya\\jeden punkt na parę};
 \node[draw=manifoldred!75!black,fill=manifoldred!7]
      (d) at (-2.55,-.90) {znoszenie uchwytów\\$W\cong N_-\times[0,1]$};
 \draw[->,thick] (a)--(b);
 \draw[->,thick] (b)--(c);
 \draw[->,thick] (c)--(d);
\end{tikzpicture}
\caption{Cztery etapy dowodu. Acykliczność
jest algebraiczna; między nią a zniesieniem
uchwytów potrzebne są ruchy geometryczne.}
\label{fig:plan-smale-19}
\end{figure}

\begin{lemat}[Brak punktów krytycznych oznacza iloczyn]
\label{lem:brak-krytycznych-hc-19}
Jeżeli funkcja Morse'a $F:W\to[0,1]$,
stała na odpowiednich brzegach,
nie ma punktów krytycznych, to
$W\cong N_-\times[0,1]$
względem $N_-$.
\end{lemat}
\begin{proof}
Wybierzmy metrykę Riemanna i pole
$X=\nabla F/|\nabla F|^2$.
Mianownik jest dodatni na całym
zwartym $W$, więc $X$ jest gładkie.
Wzdłuż jego trajektorii
$\frac{\dd}{\dd t}F(\varphi_t(x))
 =\dd F(X)=1$.
Trajektoria wychodząca z $x\in N_-$
dochodzi do poziomu $t$ po czasie $t$.
Nie urywa się wcześniej: między
poziomami leży w zwartym $W$,
a w kołnierzach pole jest poprzeczne
do brzegu. Odwrócenie czasu
przenosi każdy punkt $W$
jednoznacznie na $N_-$.
Mapa $(x,t)\mapsto\varphi_t(x)$
jest więc bijekcją, jest gładka,
a jej odwrotność
$y\mapsto(\varphi_{-F(y)}(y),F(y))$
też jest gładka.
\end{proof}

\section{Usunięcie indeksów skrajnych}
\label{sec:skrajne-hc-19}

Najpierw usuwamy uchwyty indeksów
$0,m$, potem zastępujemy indeksy
$1,m-1$ uchwytami o indeksach
bardziej środkowych. Ta druga
operacja nazywa się czasem
\emph{wymianą uchwytu}
(ang.\ \emph{handle trading}):
liczba punktów krytycznych
nie musi w tym kroku maleć.

\begin{lemat}[Znoszenie $0$- i $m$-uchwytów]
\label{lem:zero-m-hc-19}
Jeżeli oba brzegi i $W$ są spójne,
to funkcję Morse'a zgodną z brzegiem
można zmienić we wnętrzu tak, aby
nie miała punktów krytycznych
indeksu $0$ ani $m$.
\end{lemat}
\begin{proof}
Uporządkujmy uchwyty według indeksów
(Twierdzenie~\ref{thm:przestawianie-18}).
Tuż po wszystkich $1$-uchwytach
części utworzone z $N_-\times[0,1]$
i z $0$-uchwytów są wierzchołkami
grafu: każdy $1$-uchwyt jest krawędzią
łączącą dwie części albo pętlą.
Uchwyty indeksu co najmniej $2$
nie łączą różnych składowych
(ich sfery przyczepienia są spójne),
więc spójność $W$ oznacza,
że ten graf jest spójny.

Wybieramy w nim drzewo rozpinające
z korzeniem odpowiadającym $N_-$.
Dla wierzchołka reprezentującego
$0$-uchwyt wybieramy krawędź drzewa,
która łączy go w kierunku korzenia.
W pośrednim brzegu odpowiadającym
temu etapowi jedna stopa
$1$-uchwytu leży na sferze pasa
wybranego $0$-uchwytu, a druga
na innej składowej brzegu.
Sfera przyczepienia $S^0$
spotyka więc tę sferę pasa
\emph{w dokładnie jednym punkcie}.
Możemy rozdzielić wartości
krytyczne od wartości innych
uchwytów indeksu $0$ i $1$
i zastosować
Twierdzenie~\ref{thm:zniesienie-18}.
Pozostałe uchwyty są przenoszone
przez powstały pas iloczynowy.
Powtarzamy procedurę dla
wszystkich $0$-uchwytów.

W odwróconym kobordyzmie
funkcja $1-F$ zamienia indeks $m$
na $0$ oraz $m-1$ na $1$,
zgodnie z
Twierdzeniem~\ref{thm:dualny-uchwyt-18}.
Ten sam argument przy spójnym
$N_+$ usuwa wszystkie $m$-uchwyty.
\end{proof}

Sam brak indeksów $0,m$
nie wystarcza: $1$-uchwyty
zmieniają $\pi_1$ poziomic,
a dysk Whitneya dla sfer
z nimi związanych może mieć
niewłaściwy wymiar. Poniższy
krok wyjaśnia rolę prostej
spójności.

\begin{lemat}[Wymiana indeksu $1$ na $3$]
\label{lem:handel-trade-hc-19}
Jeśli $m\geq6$, $N_-$ i $W$ są
po prostu spójne, a w rozkładzie
nie ma $0$-uchwytów, to można
usunąć wszystkie $1$-uchwyty,
dopuszczając w zamian nowe
$3$-uchwyty. W odwróconym
kobordyzmie analogicznie
usuwa się indeks $m-1$
na rzecz $m-3$.
\end{lemat}
\begin{proof}
Weźmy jeden $1$-uchwyt $q$.
Po jego dołączeniu w poziomicy
jest sfera pasa $B_q\cong S^{m-2}$.
Łączymy dwie stopy uchwytu
łukiem w pozostałej części
poziomicy i zamykamy go łukiem
biegnącym przez rdzeń uchwytu.
Po niewielkim przesunięciu
otrzymujemy zanurzony okrąg
$\gamma$ przecinający $B_q$
raz i omijający sfery pasa
pozostałych $1$-uchwytów.
Jest na to miejsce, ponieważ
te sfery są rozłączne, a
wymiar poziomicy to $m-1\geq5$.

Tworzymy w małej kuli regularnego
pasa \emph{nową} parę punktów
krytycznych indeksów $2,3$.
Jest to lokalne odwrócenie
modelu~\eqref{eq:narodziny-uchwytow-18}
przy $k=2$: dla dodatniego
parametru dwa nowe punkty
są połączone jedną trajektorią.
Wybieramy kulę po starych
$2$-uchwytach i przed uchwytami
wyższych indeksów. Nowy
$2$-uchwyt ma początkowo małą
sferę przyczepienia $C\cong S^1$.

Możemy przeprowadzić $\gamma$
w górę przez poziom starych
$2$-uchwytów: okrąg i ich
sfery przyczepienia są
jednowymiarowe, więc przy
$m-1\geq5$ pozycja ogólna
pozwala je rozsunąć.
Otrzymany okrąg $\gamma'$
oraz $C$ leżą w tej samej
poziomicy po wszystkich
starych $2$-uchwytach.
Ta poziomica jest po prostu
spójna. Istotnie, jej dolna
podpoziomica ma to samo $\pi_1$
co $W$, gdyż późniejsze uchwyty
mają indeksy co najmniej $3$.
Oglądając dolną podpoziomicę
od góry, przyczepiamy do
rozważanej poziomicy jedynie
uchwyty dualne indeksu
co najmniej $m-2\geq4$,
które także nie zmieniają
$\pi_1$. Dlatego $C$ i
$\gamma'$ są homotopijne.
W rozmaitości wymiaru
$m-1\geq5$ homotopijne
zanurzone okręgi są izotopijne:
homotopię przybliża się
przez zanurzony cylinder,
bo jego ewentualne
samoprzecięcia miałyby
wymiar $2+2-(m-1)<0$.
Możemy także wybrać zanurzony
dysk ograniczony przez $\gamma'$:
nullhomotopię dysku przybliżamy
do zanurzenia z tego samego powodu.
Jego wiązka normalna jest trywialna,
bo dysk jest ściągalny. Wybieramy
obramowanie nowego $2$-uchwytu
zgodne z obramowaniem brzegu tego
dysku; w punkcie przecięcia z $B_q$
dobieramy standardowy model iloczynowy.
Izotopię z tak wybranym
obramowaniem przedłużamy na kołnierz
poziomicy i zmieniamy
pole gradientopodobne.
W nowym opisie sferą
przyczepienia nowego
$2$-uchwytu jest $\gamma'$,
a po sprowadzeniu w dół
jest nią $\gamma$.
Sfera ta przecina $B_q$
w jednym punkcie.
Twierdzenie~\ref{thm:zniesienie-18}
usuwa zatem dawny
$1$-uchwyt i nowy
$2$-uchwyt. Nowy
$3$-uchwyt zostaje.

Wykonujemy ten ruch osobno
dla kolejnych $1$-uchwytów.
Po odwróceniu całego
kobordyzmu indeksy
$1,2,3$ stają się
$m-1,m-2,m-3$;
tym samym argumentem
usuwamy $m-1$-uchwyty.
\end{proof}

\begin{figure}[htbp]
\centering
\begin{tikzpicture}[>=Latex,font=\small,
 every node/.style={align=center}]
 \node[draw=manifoldblue!75!black,fill=manifoldblue!10,
       rounded corners=3pt,minimum width=2.35cm,
       minimum height=.85cm] (a) at (-3.55,0)
       {$H_1$\\stary uchwyt};
 \node[draw=manifoldgold!85!black,fill=manifoldgold!13,
       rounded corners=3pt,minimum width=2.35cm,
       minimum height=.85cm] (b) at (0,0)
       {$H_1+H_2+H_3$\\narodziny pary};
 \node[draw=manifoldgreen!75!black,fill=manifoldgreen!11,
       rounded corners=3pt,minimum width=2.35cm,
       minimum height=.85cm] (c) at (3.55,0)
       {$H_3$\\po zniesieniu $H_1,H_2$};
 \draw[->,thick] (a)--(b);
 \draw[->,thick] (b)--(c);
\end{tikzpicture}
\caption{Wymiana $1$-uchwytu na $3$-uchwyt.
Para $H_2,H_3$ powstaje lokalnie; następnie
przesuwamy sferę przyczepienia $H_2$
i znosimy $H_1$ z $H_2$.}
\label{fig:wymiana-1-na-3-19}
\end{figure}

\begin{uwaga}[Co zostało użyte?]
\label{uwaga:wymiana-warunki-19}
Prosta spójność poziomicy służy
do przesunięcia okręgu
$\gamma'$ na $C$; bez niej
nie można uzasadnić tego kroku
samą homologią. Wymiar
co najmniej $5$ poziomicy
pozwala zastąpić homotopię
okręgów izotopią. Kwestia
normalnego obramowania
jest tu częścią wyboru
lokalnie tworzonej pary
i jej transportu przez izotopię;
to dodatkowa dana geometryczna,
której nie pokazuje schemat.
\end{uwaga}

\section{Trik Whitneya: od liczby do jednego punktu}
\label{sec:whitney-hc-19}

Po Lemacie~\ref{lem:handel-trade-hc-19}
indeksy wszystkich uchwytów leżą
między $2$ a $m-2$.
Niech $k$ będzie najmniejszym
występującym indeksem.
Ponieważ kompleks względny jest
acykliczny, jeśli pozostały uchwyty,
to $k\leq m-3$: gdyby występowały
jedynie uchwyty indeksu $m-2$,
ich wolne generatory dawałyby
niezerowe $H_{m-2}(W,N_-)$.

Ustalmy poziomicę $L$ między
uchwytami indeksów $k$ i $k+1$.
Sfery pasa $k$-uchwytów
oznaczamy $B_1,\ldots,B_r$,
a sfery przyczepienia
$(k+1)$-uchwytów
$A_1,\ldots,A_s$.
Mają w $L^{m-1}$ wymiary
\[
 \dim A_j=k,\qquad
 \dim B_i=m-k-1,\qquad
 \dim A_j+\dim B_i=m-1.
\]
Po transwersalnej perturbacji
przecinają się w skończenie
wielu punktach. Współczynnik
macierzy $\partial_{k+1}$
jest \emph{sumą ich znaków},
jak w~\eqref{eq:macierz-uchwyty-18}.

\begin{lemat}[Prosta spójność poziomicy i jej dopełnienia]
\label{lem:komplement-whitney-19}
W sytuacji z poprzedniego akapitu
załóżmy, że $N_-$ i $N_+$ są
po prostu spójne, a $2\leq k\leq m-3$.
Wtedy $L$ oraz dopełnienie
\[
 L\setminus
   \left(\bigcup_i B_i\ \cup\ \bigcup_j A_j\right)
\]
są po prostu spójne.
\end{lemat}
\begin{proof}
Oznaczmy przez $W_-$ część
kobordyzmu od $N_-$ do $L$.
Ponieważ przed $L$ przyczepiano
tylko uchwyty indeksu $k$,
a $k\geq2$, te uchwyty nie mogą
\emph{utworzyć} nowej grupy
fundamentalnej ze spójnego,
po prostu spójnego $N_-$:
$2$-komórka może zabić pętlę,
a komórki wyższego wymiaru
nie zmieniają $\pi_1$.
Stąd $\pi_1(W_-)=0$.
Odwróćmy fragment $W_-$.
Jest zbudowany z $L$
przez dołączanie uchwytów
dualnych indeksu $m-k\geq3$.
One nie zmieniają $\pi_1$,
zatem $\pi_1(L)=0$.

Jeżeli $2<k<m-3$, to
$\operatorname{codim}_L B_i=k\geq3$
i $\operatorname{codim}_L A_j
=m-k-1\geq3$.
Usuwanie skończonej sumy
takich podrozmaitości nie
zmienia $\pi_1$.
Aby to zobaczyć, przesuwamy
pętlę na mocy transwersalności
poza te zbiory.
Jej homotopia wypełniająca
ma wymiar $2$ i również
może ich uniknąć, ponieważ
$2+\dim B_i<\dim L$
i $2+\dim A_j<\dim L$.
Dlatego dopełnienie jest
po prostu spójne.

Trzeba osobno sprawdzić
$k=2$, gdy $B_i$ ma
kowymiar $2$ w $L$.
Usuwamy najpierw wszystkie
$B_i$. W modelu $2$-uchwytu
nowa część brzegu ma postać
$D^2\times S^{m-3}$, a sfera pasa
to $\{0\}\times S^{m-3}$.
Po jej usunięciu dyskowy czynnik
odkształca się do $S^1$.
Na starym brzegu analogiczny obszar
to otoczenie okręgu przyczepienia
$S^1\times D^{m-2}$ z usuniętym
rdzeniem. Sklejone kołnierze mają
ten sam typ homotopijny. Regularny
przepływ identyfikuje więc grupę
$\pi_1(L\setminus\bigcup B_i)$
z grupą $N_-$ po usunięciu
okręgów przyczepienia
odpowiednich $2$-uchwytów.
W $N_-^{m-1}$ mają one
kowymiar $m-2\geq4$;
z poprzedniego argumentu
ich usunięcie nie zmienia
$\pi_1(N_-)=0$.
Następnie usuwamy $A_j$,
których kowymiar wynosi
$m-3\geq3$, więc $\pi_1$
pozostaje zerowa.
Dla $k=m-3$ odwracamy
kobordyzm i stosujemy
ten sam argument do
drugiego brzegu.
\end{proof}

\begin{lemat}[Trik Whitneya dla rozważanych sfer]
\label{lem:whitney-hc-19}
W warunkach
Lematu~\ref{lem:komplement-whitney-19}
niech $A_j$ i $B_i$
przecinają się w dwóch
punktach $u,v$ o przeciwnych
znakach. Można zmienić
pole gradientopodobne
w małym pasie wokół $L$,
tak aby te dwa punkty
zniknęły, bez zmiany innych
punktów przecięcia
którejkolwiek pary sfer.
\end{lemat}
\begin{proof}
\emph{Krok 1: brzeg dysku.}
Wybieramy łuk $\alpha\subset A_j$
od $u$ do $v$ i łuk
$\beta\subset B_i$ w przeciwnym
kierunku. Ich wnętrza
omijają pozostałe punkty
przecięcia: obie sfery mają
wymiar co najmniej $2$,
a usuwamy zeń tylko
skończony zbiór.
Suma $\alpha\cup\beta$
jest zamkniętą krzywą
o dwóch narożach.
W małym otoczeniu tej krzywej
odsuwamy jej równoległą
kopię od obu rodzin sfer.
Ta kopia leży w dopełnieniu
z poprzedniego lematu,
więc jest tam homotopijnie
trywialna.

\emph{Krok 2: wypełnienie.}
Wybierzmy mapę dysku
rozpinającą odsuniętą pętlę
w tym dopełnieniu.
Transwersalność względna
pozwala uczynić ją
zanurzeniem, bo oczekiwany
wymiar samoprzecięć
dwóch fragmentów
dwuwymiarowego dysku
wynosi $2+2-(m-1)<0$.
Dołączając wąski kołnierz
od równoległej kopii do
$\alpha\cup\beta$
i wygładzając naroża,
otrzymujemy zanurzony
dysk $D$ o brzegu leżącym
na $A_j\cup B_i$,
którego wnętrze omija
\emph{wszystkie} $A$ i $B$.
Jest to dysk Whitneya.

\emph{Krok 3: obramowanie.}
Samo zanurzenie dysku $D$ nie
wystarcza do wykonania ruchu.
Wzdłuż $\alpha$ i $\beta$ obie sfery
wyznaczają kierunki, w których
należy przesunąć ich pasy. Potrzebne
jest przedłużenie tych danych do
\emph{obramowanego dysku Whitneya}.
Zgodność znaków na końcach łuków
jest konieczna, lecz sama nie
dowodzi istnienia takiego przedłużenia:
trzeba także skontrolować skręcenie
normalne wzdłuż całego brzegu.
Stosujemy tu standardowy \emph{obramowany
trik Whitneya} dla dwóch sfer wymiarów
co najmniej $2$ w rozmaitości wymiaru
co najmniej $5$. Jego lokalna część
pozwala skorygować obramowanie dysku
przy brzegu; w skrajnych kowymiarach
potrzebny jest dodatkowo warunek o
$\pi_1$ dopełnienia, sprawdzony w
Lemacie~\ref{lem:komplement-whitney-19}.
Nie wynika to wyłącznie z prostego
rachunku $\pi_1$ przestrzeni
Grassmanna. Pełna wersja lematu
geometrycznego jest użyta również w
\href{https://him-lueck.uni-bonn.de/data/Chapter_s-Cobordism_theorem.pdf}
{dowodzie Lücka i Macko, Lemat 2.35}.
Otrzymujemy obramowane
otoczenie
\[
 D^2\times D^{k-1}
       \times D^{m-k-2}
\]
o wymiarze $m-1$,
w którym $A_j$ i $B_i$
są dwiema standardowymi
poprzecznymi rodzinami pasów.

\emph{Krok 4: ruch.}
W tym otoczeniu przesuwamy
pas $A_j$ przez $D^2$,
utrzymując go nieruchomo
blisko końców i na zewnątrz
otoczenia. W płaskim
przekroju wygląda to jak
usunięcie dwóch przecięć
z Rysunku~\ref{fig:whitney-hc-19}.
Iloczynowa postać normalnych
kierunków zapewnia,
że ten sam ruch jest izotopią
w pełnym wymiarze $m-1$.
Ponieważ wnętrze dysku
nie spotyka pozostałych sfer,
nie powstają dodatkowe
punkty przecięcia.
Przedłużamy izotopię
na regularny pas przepływu
i modyfikujemy tam pole
gradientopodobne; funkcja
i jej punkty krytyczne
pozostają bez zmian.
\end{proof}

\begin{figure}[htbp]
\centering
\begin{tikzpicture}[>=Latex,font=\small,line cap=round,scale=1.0]
 \node[font=\bfseries] at (-2.4,1.18) {przed};
 \node[font=\bfseries] at (2.55,1.18) {po};
 \begin{scope}[shift={(-4.7,0)}]
  \fill[manifoldgold!22]
    (.85,0) .. controls (1.20,.78) and (2.20,.78)
       .. (2.55,0) -- cycle;
  \draw[manifoldblue!83!black,very thick]
    (.05,0)--(3.40,0);
  \draw[manifoldred!82!black,very thick]
    (.35,-.65) .. controls (.50,-.35) and (.64,-.16)
     .. (.85,0) .. controls (1.20,.78) and (2.20,.78)
     .. (2.55,0) .. controls (2.77,-.16)
         and (2.90,-.35) .. (3.05,-.65);
  \fill[manifoldgold!83!black] (.85,0) circle (2.4pt);
  \fill[manifoldgold!83!black] (2.55,0) circle (2.4pt);
  \node[anchor=south] at (.85,.10) {$+$};
  \node[anchor=south] at (2.55,.10) {$-$};
  \node[manifoldgold!75!black] at (1.70,.40) {$D$};
  \node[manifoldblue!80!black,anchor=north] at (.35,-.04) {$B$};
  \node[manifoldred!80!black,anchor=west] at (3.05,-.61) {$A$};
 \end{scope}
 \draw[->,gray!75,thick] (-.45,-.04)--(.42,-.04);
 \begin{scope}[shift={(1.0,0)}]
  \draw[manifoldblue!83!black,very thick]
    (.05,0)--(3.40,0);
  \draw[manifoldred!82!black,very thick]
    (.35,-.65) .. controls (1.15,-.18)
           and (2.25,-.18) .. (3.05,-.65);
  \node[manifoldblue!80!black,anchor=north] at (.35,-.04) {$B$};
  \node[manifoldred!80!black,anchor=west] at (3.05,-.61) {$A$};
 \end{scope}
\end{tikzpicture}
\caption{Dwuwymiarowy przekrój triku Whitneya:
punkty o przeciwnych znakach są końcami
brzegu złotego dysku $D$. Przesunięcie
czerwonej sfery usuwa oba punkty.
Sam dowód działa w poziomicy
wymiaru co najmniej $5$;
rysunek nie przedstawia jej
pozostałych wymiarów.}
\label{fig:whitney-hc-19}
\end{figure}

\begin{przyklad}[Dlaczego $\pm1$ to za mało bez triku?]
Niech para sfer ma trzy punkty
przecięcia o znakach $+,+,-$.
Ich liczba algebraiczna to $1$,
lecz geometryczna liczba wynosi $3$.
Trik usuwa parę $+,-$ i pozostawia
jeden punkt $+$. Dopiero teraz
Twierdzenie~\ref{thm:zniesienie-18}
pozwala znieść uchwyty.
\end{przyklad}

\begin{uwaga}[Znaczenie założeń wymiarowych]
Dla wymiaru poziomicy $4$
wypełniający dysk może sam siebie
przecinać, ponieważ $2+2-4=0$.
Dlatego argument z Kroku 2
nie dowodzi twierdzenia o
gładkich $5$-wymiarowych
$h$-kobordyzmach.
Nie zakładamy też, że
zerowa liczba przecięcia
zawsze daje dysk Whitneya:
potrzebna była prosta
spójność \emph{dopełnienia}
całej rodziny sfer.
\end{uwaga}

\section{Od macierzy do geometrycznych par}
\label{sec:macierze-hc-19}

Po zniesieniu indeksów skrajnych
najmniejszy indeks oznaczmy $k$.
Brak $C_{k-1}$ oraz
$H_k(C_\bullet)=0$ dają
$\operatorname{im}\partial_{k+1}
=\ker\partial_k=C_k$.
Macierz
\[
 \partial_{k+1}:\Z^s\longrightarrow\Z^r,
 \qquad r=\#\{k\text{-uchwytów}\},
\]
jest więc \emph{surjektywna}.
Nad $\Z$ wolno wykonywać
wyłącznie operacje z
całkowitymi współczynnikami;
dzielenie przez liczbę $2$
zmieniłoby problem geometryczny.

\begin{lemat}[Całkowita redukcja przez kolumny]
\label{lem:macierz-surj-hc-19}
Niech $A:\Z^s\to\Z^r$
będzie surjekcją.
Permutacje kolumn, zmiana
znaku kolumny i dodawanie
całkowitej wielokrotności
jednej kolumny do drugiej
przeprowadzają jej macierz
do postaci
\[
       \begin{pmatrix}I_r&0\end{pmatrix}.
\]
Nie trzeba zmieniać
wybranej bazy $\Z^r$.
\end{lemat}
\begin{proof}
Dowodzimy przez indukcję po $r$.
Dla $r=0$ nie ma czego robić.
Pierwszy wiersz macierzy
surjektywnej ma wpisy
o największym wspólnym
dzielniku $1$: pierwszy
współrzędny obraz całej
$\Z^s$ jest równy $\Z$.
Algorytm Euklidesa można
wykonać operacjami na
\emph{kolumnach}.
Jeśli w dwóch miejscach
pierwszego wiersza stoją
$a,b$ z $b\ne0$, dzielimy
$a=qb+t$, $|t|<|b|$,
i zastępujemy odpowiednią
kolumnę przez nią minus
$q$ razy drugą. Powtarzając
otrzymujemy wpis $\pm1$.
Permutujemy go do pierwszej
kolumny i zmieniamy znak,
aby uzyskać $1$.
Odejmując wielokrotności
pierwszej kolumny od
pozostałych, zerujemy
resztę pierwszego wiersza.

Zapiszmy teraz
\[
 A=\begin{pmatrix}
       1&0\\
       b&A'
    \end{pmatrix}.
\]
Macierz $A'$ z $r-1$
wierszami jest surjektywna.
Istotnie, dla dowolnego
$y\in\Z^{r-1}$ wektor
$(0,y)$ należy do obrazu $A$.
Jego przeciwobraz ma
pierwszą współrzędną
równą $0$, gdyż pierwszy
wiersz to $(1,0,\ldots,0)$;
zatem $y$ leży w obrazie
$A'$. Z założenia
indukcyjnego operacje
na kolumnach poza pierwszą
zamieniają $A'$ w
$(I_{r-1}\ 0)$ i nie
zmieniają pierwszego wiersza.
Na koniec odejmujemy
od pierwszej kolumny
odpowiednią kombinację
kolumn z tym $I_{r-1}$,
aby wyzerować $b$.
Dostajemy $(I_r\ 0)$.
Każda użyta operacja
ma całkowitą odwrotność.
\end{proof}

\begin{przyklad}[Dwa przesunięcia, bez dzielenia]
Niech w wybranych bazach
\[
 [\partial_{k+1}]
 =\begin{pmatrix}2&1\\1&1\end{pmatrix},
 \qquad
 \partial e_1=2f_1+f_2,\quad
 \partial e_2=f_1+f_2.
\]
Wybieramy
$e'_1=e_1-e_2$, zatem
$\partial e'_1=f_1$.
Następnie
$e'_2=e_2-e'_1=2e_2-e_1$,
więc $\partial e'_2=f_2$.
Obie zmiany są przesunięciami
uchwytów indeksu $k+1$.
Macierz stała się $I_2$,
ale dopiero trik Whitneya
usunie ewentualne
dodatkowe pary
geometrycznych przecięć.
\end{przyklad}

\begin{stwierdzenie}[Geometryczna realizacja redukcji]
\label{prop:realizacja-hc-19}
Po wykonaniu zmian z
Lematu~\ref{lem:macierz-surj-hc-19}
można doprowadzić do sytuacji,
w której sfery
$A_1,\ldots,A_s$ i
$B_1,\ldots,B_r$
spotykają się geometrycznie
dokładnie tak:
\[
 \#(A_j\cap B_i)=
 \begin{cases}
  1,& i=j\leq r,\\
  0,& \text{w przeciwnym razie}.
 \end{cases}
\]
Pozostałe części funkcji Morse'a
i wcześniej zbudowany brzeg
wejściowy można zachować.
\end{stwierdzenie}
\begin{proof}
Zmiana kolejności kolumn
to zmiana numeracji
$(k+1)$-uchwytów,
a zmiana znaku kolumny
to odwrócenie orientacji
rdzenia. Operację
$e_p\mapsto e_p\pm e_q$
realizuje przesunięcie
uchwytu z
Twierdzenia~\ref{prop:slide-18}.
Wielokrotność całkowitą
otrzymujemy przez
skończone powtórzenie
przesunięcia.
Wszystkie te ruchy
zachowują dyfeomorfizm
kobordyzmu względem
wcześniejszego brzegu.
Po nich algebraiczna
liczba przecięć
$I(A_j,B_i)$ jest
$\delta_{ij}$ dla $i\leq r$,
a dla pozostałych par
wynosi zero.

Każda para punktów
o przeciwnych znakach
na ustalonej parze sfer
może być usunięta
Lematem~\ref{lem:whitney-hc-19},
bez wytworzenia przecięć
z innymi sferami.
Dla $I=0$ wszystkie
punkty znikają parami.
Dla $I=1$ po ich
usunięciu zostaje
jeden punkt dodatni.
Liczba przecięć jest
skończona, więc
proces kończy się
po skończenie wielu
ruchach. Otrzymujemy
zapowiedzianą postać
geometryczną.
\end{proof}

\section{Dowód twierdzenia i jego granice}
\label{sec:dowod-smale-19}

\begin{proof}[Dowód Twierdzenia~\ref{thm:smale-hc-19}]
\emph{Krok 1: wybór rozkładu.}
Bierzemy funkcję Morse'a
zgodną z kołnierzami,
której istnienie dało
Stwierdzenie~\ref{prop:morse-kobordyzm-18}.
Rozkład na uchwyty
i jego kompleks względny
pochodzą z
Twierdzeń~\ref{thm:rozklad-uchwyty-18}
i~\ref{thm:kompleks-uchwytow-18}.
Stwierdzenie~\ref{prop:acykliczny-hc-19}
daje acykliczność.
Ustawiamy indeksy
niemalejąco, korzystając
z~\ref{thm:przestawianie-18}.

\emph{Krok 2: usuwamy indeksy skrajne.}
Ponieważ oba brzegi są
spójne, Lematem~\ref{lem:zero-m-hc-19}
usuwamy indeksy $0,m$.
Założenia o prostej
spójności i $m\geq6$
pozwalają zastosować
Lemat~\ref{lem:handel-trade-hc-19}
oraz jego wersję dualną:
odtąd występują tylko
indeksy $2,\ldots,m-2$.
Wszystkie te zmiany
modyfikują funkcję
i pole, ale nie $W$
ani jego brzegu wejściowego.

\emph{Krok 3: usuwamy najniższy indeks.}
Jeżeli są jeszcze punkty
krytyczne, niech $k$
będzie najmniejszym
indeksem. Acykliczność
wyklucza $k=m-2$;
zatem $2\leq k\leq m-3$.
Nie ma uchwytów indeksu
$k-1$, więc
$\partial_{k+1}:C_{k+1}\to C_k$
jest surjekcją.
Lemat~\ref{lem:macierz-surj-hc-19}
i Stwierdzenie~\ref{prop:realizacja-hc-19}
pozwalają ułożyć
$k$- i $(k+1)$-uchwyty
w geometryczne pary,
w których każda
sfera przyczepienia
spotyka odpowiednią
sferę pasa dokładnie raz
i omija pozostałe.

Rozdzielamy wartości
krytyczne pierwszej pary
od wartości innych
punktów tego samego
indeksu. Stosujemy
Twierdzenie~\ref{thm:zniesienie-18}
i usuwamy tę parę.
Ponieważ poprzedni
układ sfer nie miał
innych przecięć,
powtarzamy operację
dla pozostałych par.
Wynikowy kompleks
pozostaje acykliczny,
gdyż nadal oblicza
$H_*(W,N_-)=0$,
a liczba punktów
indeksu $k$ spadła
do zera.

\emph{Krok 4: kończymy indukcję.}
Przechodzimy do następnego
najmniejszego indeksu.
W każdym cyklu znika
co najmniej jedna
para punktów krytycznych,
a nie wracamy już do
niższych indeksów.
Ponieważ rozkład ma
skończenie wiele uchwytów,
po skończenie wielu
krokach nie ma punktów
krytycznych. Z
Lematu~\ref{lem:brak-krytycznych-hc-19}
otrzymujemy dyfeomorfizm
$\Phi:N_-\times[0,1]\to W$
identyczny na wejściowym
brzegu. Ograniczenie
$\Phi|_{N_-\times\{1\}}$
jest dyfeomorfizmem
$N_-\to N_+$.
\end{proof}

\begin{uwaga}[Dlaczego dowód nie kończy się na macierzy?]
Algebra zapewnia
$[\partial_{k+1}]=(I\ 0)$.
Taka macierz mówi o
\emph{podpisanych sumach},
podczas gdy twierdzenie
o znoszeniu potrzebuje
\emph{jednego punktu}.
Lemat~\ref{lem:komplement-whitney-19}
zapewnia miejsce
na dyski Whitneya,
a Lemat~\ref{lem:whitney-hc-19}
usuwa pary zbędnych
punktów. To tu
prosta spójność i
wysoki wymiar
wchodzą do argumentu
w sposób zasadniczy.
\end{uwaga}

\begin{przyklad}[Iloczyn z niepotrzebną parą uchwytów]
Zacznijmy od
$N\times[0,1]$ dla
po prostu spójnej
zamkniętej $N^{m-1}$.
Lokalna zmiana
funkcji bez ruszania
brzegu może utworzyć
parę punktów
indeksów $2,3$,
jak w modelu
\eqref{eq:narodziny-uchwytow-18}.
Jej kompleks względny
ma fragment
$\Z\xrightarrow{\ \pm1\ }\Z$.
Przestrzeń była i jest
iloczynem; zapis
Morse'a tylko
chwilowo opisuje ją
dwoma uchwytami.
Znoszenie pary
przywraca funkcję
bez punktów krytycznych.
\end{przyklad}

\begin{wniosek}[Sfery homotopijne w wysokim wymiarze]
\label{wn:homotopijna-sfera-19}
Jeśli zamknięta gładka
$\Sigma^n$, $n\geq6$,
ma typ homotopijny
$S^n$, to jest
homeomorficzna z $S^n$.
Dokładniej: jest sumą
dwóch gładkich dysków
$D^n$ sklejonych
pewnym dyfeomorfizmem
ich brzegów.
\end{wniosek}
\begin{proof}
Wybieramy rozłączne
zanurzone dyski
$B_-,B_+\subset\Sigma$
i usuwamy ich wnętrza:
\[
 W=\Sigma\setminus
   \bigl(\operatorname{int}B_-
          \cup\operatorname{int}B_+\bigr).
\]
Brzegi $W$ są sferami
$S^{n-1}$.
Ponieważ $n\geq6$,
przyklejenie dwóch
$n$-dysków nie zmienia
grupy fundamentalnej;
z twierdzenia van Kampena
$\pi_1(W)=\pi_1(\Sigma)=0$.
Ciąg Mayera--Vietorisa
dla pokrycia $\Sigma$
przez $W$ i oba dyski
pokazuje, że
$\widetilde H_j(W)=0$
dla $j\ne n-1$, a
$H_{n-1}(W)\cong\Z$.
Dokładniej, mapa
\[
 H_n(\Sigma)\cong\Z
 \longrightarrow
 H_{n-1}(S^{n-1}\sqcup S^{n-1})
      \cong\Z^2
\]
wysyła klasę fundamentalną
na parę klas brzegowych
o znakach przeciwnych.
Wynikowy iloraz jest
$\Z$, a każda z dwóch
sfer brzegowych
odwzorowuje się na
jego generator.
Obie inkluzje
$S^{n-1}\hookrightarrow W$
są więc izomorfizmami
homologii. Ponieważ
są to po prostu spójne
kompleksy typu CW,
twierdzenie Whiteheada
w wersji homologicznej
(por.\ Rozdział~\ref{ch:wyzsza-homotopia})
czyni je równoważnościami
homotopijnymi.
Stąd $W$ jest
$h$-kobordyzmem.
Twierdzenie~\ref{thm:smale-hc-19}
daje $W\cong S^{n-1}\times[0,1]$.
Po doklejeniu $B_-$ mamy
dysk; $B_+$ doklejamy
do jego brzegu pewnym
dyfeomorfizmem
$\psi:S^{n-1}\to S^{n-1}$.
Topologicznie $\psi$
przedłuża się na dysk
wzorem radialnym
$r u\mapsto r\psi(u)$.
Po takim przedłużeniu
dwa sklejone dyski
dają zwykłą sferę
topologiczną.
\end{proof}

\begin{uwaga}[Topologia i gładkość sfery]
Ostatni wniosek orzeka
homeomorfizm ze standardową
sferą, niekoniecznie
dyfeomorfizm. Przedłużenie
radialne jest ciągłe,
lecz zwykle nie jest
gładkie w środku dysku.
Dlatego argument pozostawia
miejsce dla egzotycznych
struktur gładkich
na sferach.
\end{uwaga}

\section{Co przygotowaliśmy do dalszej teorii?}
\label{sec:dalej-po-smale-19}

W dowodzie znikały kolejno
cztery trudności:
spójność składowych przy
$0$- i $m$-uchwytach,
pętle związane z
$1$- i $(m-1)$-uchwytami,
niejednostkowa macierz
$\partial_{k+1}$ oraz
zbędne geometryczne
przecięcia sfer.
Rozdział~\ref{ch:rachunek-uchwytow}
wyjaśnił pojedynczy ruch;
niniejszy rozdział złożył
ruchy w dowód globalny.

Gdy $\pi_1(W)$ nie jest
trywialna, liczby całkowite
w macierzy przecięć
zastępują współczynniki
w pierścieniu
$\Z[\pi_1(W)]$.
Samo zerowanie homologii
nie wystarcza wtedy
do uproszczenia macierzy
przez dozwolone ruchy.
Dodatkową przeszkodą
jest \emph{torsja Whiteheada},
której definicja oraz
twierdzenie o
$s$-kobordyzmie będą
naturalnym następnym
krokiem. Jest to inne
pojęcie niż torsja
elementu modułu
z Rozdziału~\ref{ch:algebra-abstrakcyjna}.
Z kolei ślad chirurgii
z Twierdzenia~\ref{thm:slad-chirurgii-18}
pozwala później
przechodzić od zmian
kobordyzmu do
kontrolowanych operacji
na jego poziomicach.

\par\smallskip
{\footnotesize\noindent\emph{Dalsza lektura:}
\href{https://www.math.columbia.edu/~jmorgan/Lecture_III_H-cobordism_Theorem.pdf}
{J.~Morgan, \emph{Lecture III: h-cobordism Theorem}},
\href{https://www.math.columbia.edu/~jmorgan/Lecture_IIIA_hcobordism_Contd.pdf}
{J.~Morgan, \emph{Lecture IIIA: cancellation and Whitney trick}},
\href{https://www.cambridge.org/core/books/differential-topology/theory-of-handle-decompositions/BA54EF1BD9B755DFA2B363176B0CA2F7}
{C.\,T.\,C.~Wall, \emph{Differential Topology}, rozdział o uchwytach}.\par}
